\documentclass[12pt, twoside, openany]{book}

\usepackage{mathptmx} % Contiene una fuente similar a Times New Roman

\usepackage[spanish, es-tabla]{babel} % Permite escritura en castellano
\usepackage[utf8]{inputenc} % Permite utilizar caracteres UTF8

\usepackage{graphicx} % Para la inclusión de gráficos e imágenes
\graphicspath{ {images/} } % Ruta para buscar las imágenes

\usepackage[a4paper,top=30mm,left=30mm,right=25mm,bottom=25mm,headheight=20mm]{geometry} % Configuración de los margenes de la página

% Paquetes para que funcione el formato. Si se tiene dudas con el funcionamiento de alguno o se quiere configurar de manera diferente, se recomienda leer la ayuda de LaTeX sobre estos paquetes
\usepackage{titlesec}
\usepackage{setspace}
\usepackage{ragged2e}
\usepackage{fancyhdr}
\usepackage{lastpage}
\usepackage{stackengine}
\usepackage{array}
\usepackage{url}
\usepackage{float}

% Paquete para que las referencias funcionen, y permite introducir links
\usepackage{hyperref}
% Paquete para trabajar con colores (fondo de celdas, color del texto...)
\usepackage{xcolor}

% Paquete para incluir páginas de otro PDF directamente
\usepackage{pdfpages}

% Paquete para generación de texto
\usepackage{lipsum}

% Se carga el fichero de configuración (se ha separado para facilitar la comprensión de este ejemplo)
\input{Configuracion.tex}

\begin{document}

\rmfamily % Fuente tipo Romana

% Portada de la memoria
\input{Portada.tex}

% Índice de contenido
\addcontentsline{toc}{chapter}{Índice de contenido} % Añade la referencia al índice de contenido
\tableofcontents
\newpage

% Índice de figuras
\addcontentsline{toc}{chapter}{Índice de figuras}  % Añade la referencia al índice de contenido
\listoffigures
\newpage

% Índice de tablas
\addcontentsline{toc}{chapter}{Índice de tablas} % Añade la referencia al índice de contenido
\listoftables

\justify % Texto justificado

\setlength{\parskip}{\baselineskip} % Separación entre párrafos de 1 linea
\onehalfspacing % Interlineado de 1,5 

%% Aunque en el ejemplo he creado varios capítulos en el fichero 01_Introduccion.tex, suelo preferir crear varios ficheros, uno para cada capitulo. Considero que reduciendo el tamaño de los ficheros se puede trabajar más cómodamente
\input{01_Introduccion.tex} % Incluye el contenido del fichero 01_Introduccion.tex

%% Esto incluirá la bibliografía correctamente en nuestro trabajo
\newpage % En una nueva página
\addcontentsline{toc}{chapter}{Bibliografía} % Añade la referencia al índice de contenido

\bibliographystyle{ieeetr} % Define el estilo de la bibliografía
\bibliography{bibliografia} % Indica el archivo que contiene la colección de citas
% Cualquier otro aspecto que se quiera modificar con la bibliografía puede ser útil consultar la ayuda de LaTeX

\end{document}